% !TEX root = ../jlt.tex
\section{From Featherweight Jlt to System~F}
\label{sec:translation}
% Establishes: what a reader must believe about the System F reading, by one
% example and the rules it uses.  The full calculus (the appendices)
% is the warranty; a reader who stops here knows the claims and their cost.

This section is about one map, the translation $\llbracket \cdot \rrbracket$ from Featherweight Jlt into System~F with records and the iterators of inductive types, defined in full in Appendix~\ref{app:fjlt:translation}.
It is total on FJlt, and Section~\ref{sec:guarantees} says what the image guarantees; it is partial on \Cpp{}, and Section~\ref{sec:translation:adequacy} is the argument for when a \Cpp{} program is an FJlt program.
Six of its rules carry the listings' structure, and Table~\ref{tab:rules} states them.
Conditionals, conversions and the remaining operators are routine, and Appendix~\ref{app:fjlt:translation} gives them all.

\begin{table}[ht]
\centering
\small
\begin{tabularx}{\linewidth}{@{}>{\raggedright\arraybackslash}p{0.36\linewidth}>{\raggedright\arraybackslash}X>{\raggedright\arraybackslash}p{0.17\linewidth}@{}}
\toprule
FJlt & System~F & in \\
\midrule
\texttt{concept C = R̄} & $\mathit{Dict}_C\,\overline{\alpha} = \{\,\overline{r : \llbracket R \rrbracket}\,\}$ & Listing~\ref{lst:rule2} \\
\texttt{constexpr τ f(τ̄) = delete} & nothing & Listing~\ref{lst:rule2} \\
\texttt{template <>}\newline\texttt{constexpr τ f<τ̄>(x̄) \{ e \}} & $\lambda\overline{x}.\,\llbracket e \rrbracket$ & Listing~\ref{lst:rule2} \\
\texttt{template <ᾱ> requires Q̄}\newline\texttt{constexpr τ f(x̄) \{ e \}} & $\Lambda\overline{\alpha}.\,\lambda\overline{q : \mathit{Dict}_Q}.\,\lambda\overline{x}.\,\llbracket e \rrbracket$ & Listing~\ref{lst:rule2} \\
\texttt{f<τ̄>(ē)} & $\llbracket f \rrbracket\,\overline{\llbracket \tau \rrbracket}\,\overline{d}\,\overline{\llbracket e \rrbracket}$ & Listing~\ref{lst:rule2} \\
\texttt{fold\_left(r, a, f)}, $r$ a chain & $\pi_2\,(\mathit{iter}\,|r|\,\mathit{step}\,(\llbracket r \rrbracket, \llbracket a \rrbracket))$, $\mathit{step}$ the chain's $\mathit{next}$ beside $\llbracket f \rrbracket$ at its $\mathit{head}$ & Listing~\ref{lst:rule3} \\
\bottomrule
\end{tabularx}
\caption{Six translation rules exercised by the listings; a chain is iterated for its length, its position carried beside the accumulator (Figure~\ref{fig:structures}).
The whole translation is Figure~\ref{fig:fjlt-translation} in Appendix~\ref{app:fjlt:translation}.}
\label{tab:rules}
\end{table}

A law is a field of type \texttt{Bool}, and a shape built from shapes is a record holding their records, which is how a superclass is passed.
The whole is dictionary passing, Wadler and Blott's translation of type classes, here into System~F, the record being the dictionary~\cite{wadler1989adhoc}.
Table~\ref{tab:dictionary} names the \Cpp{} devices it meets.
\cppinline{Op} is a type with exactly one value, the closure \cppinline{Op{}}, so under translation the type \emph{is} that closure, and \cppinline{Op{}(a, b)} is \texttt{d.op a b}.

\begin{table}[ht]
\centering
\small
\begin{tabularx}{\linewidth}{@{}XXX@{}}
\toprule
\Cpp{23} & type-class reading & under translation \\
\midrule
concept & class & record type \\
requirement met by the type's own operation & an instance nobody wrote & field filled by a primitive \\
explicit specialisation of a deleted function template & instance & field supplied by a value \\
constrained overload & instance with a context & function from records to a record \\
\texttt{requires} clause on a function & class constraint & one more parameter \\
\bottomrule
\end{tabularx}
\caption{Dictionary passing, met in \Cpp{}.
The second row has no counterpart in a nominal type-class system: concepts are structural, so a type that has the operation satisfies the shape without a declaration.}
\label{tab:dictionary}
\end{table}

\subsection{Instances with contexts}
\label{sec:translation:contexts}

An instance may itself require instances: the unit of a pair is the pair of units, provided both components have one.
Listing~\ref{lst:rule2}'s instance for the unsigned integrals is the simplest case, one context on one carrier.
In \Cpp{} that is a generic function with a \cppinline{requires} clause, and the clause is the context.
Under translation it is a function from the components' records to the composite's:
\begin{minted}{text}
pair_unit  =  Λ A. Λ B. λ da : MonoidDict A OpA. λ db : MonoidDict B OpB.
                (da.unit, db.unit)
\end{minted}
The compiler applies such functions outside in, as often as the composite nests.
The recursion ends, because a context on a structured head mentions only proper parts of its type arguments, and a context mentions only shapes declared before it.\footnote{The termination conditions on instance contexts are studied by Sulzmann et al.~\cite{sulzmann2007understanding}.
GHC calls them the Paterson conditions.
FJlt's well-formedness requires them (Appendix~\ref{app:fjlt:typing}).}

